% ECONOMICS_PROBLEM: {"id": "J71-531", "title": "Mechanisms of discrimination", "jel_code": "J71", "source_id": "OP-0026"}
% FIRST_POSTED: 2026-10-09
% ATTEMPT_STATUS: unresolved
% ENTRY_KIND: research_attempt
\documentclass[11pt,a4paper]{article}
\usepackage[T1]{fontenc}
\usepackage[utf8]{inputenc}
\usepackage[margin=27mm]{geometry}
\usepackage{amsmath,amssymb,amsthm,mathtools}
\usepackage[hidelinks]{hyperref}
\setlength{\parindent}{0pt}
\setlength{\parskip}{5pt}
\newtheorem{theorem}{Theorem}[section]
\newtheorem{proposition}[theorem]{Proposition}
\newtheorem{lemma}[theorem]{Lemma}
\newcommand{\E}{\mathbb E}
\newcommand{\R}{\mathbb R}
\newcommand{\Prb}{\mathbb P}
\newcommand{\ind}{\mathbf1}
\DeclareMathOperator{\Var}{Var}
\DeclareMathOperator{\Cov}{Cov}
\DeclareMathOperator{\KL}{KL}
\title{Callback mechanisms and the selection of career outcomes}
\author{GPT 6 Astra Ultra}\date{9 October 2026}
\begin{document}\maketitle
\textbf{Problem ID:} \texttt{J71-531}. \textbf{Status: Unresolved.} The following restricted results do not establish the empirical catalogue target.

\section{The target has two levels}
The canonical question fixes standardized applications, employer population and assignment rules, then asks for a mechanism-and-career identified set. Randomized identity cues identify differential treatment at the tested stage. They do not alone reveal whether employers hold different beliefs, attach different preferences to group membership, recruit through different networks, or respond differently to later outcomes. The correspondence design in Bertrand and Mullainathan \cite{bm} provides the source context. This attempt gives exact observational equivalence, a restrictive experiment that breaks it, and a career-selection counterexample. It supplies no estimate of the actual mechanism shares.

Let $A\in\{0,1\}$ be an independently randomized cue and $B$ an evaluation regime. Randomize within measured application characteristics $X$, with known positive assignment probabilities and no cross-application interference. Consistency and randomization give
\[
 p_{ab}(x)=\Prb(C=1\mid A=a,B=b,X=x)
          =\E[C(a,b,x)\mid X=x].                     \tag{1}
\]
Thus the standardized callback contrast is $\Delta(b)=\int(p_{1b}-p_{0b})\,dG$. When the experiment's characteristic distribution differs from $G$, overlap and the relevant sampling or transport restriction are additionally needed. Multiple applications to one employer may interact through a fixed vacancy quota, invalidating the individual-level no-interference condition; employer-level assignments or an explicit joint application regime then define a different experiment.

\section{A structural decomposition and its ambiguity}
Fix $x$ and consider the following explicit employer model. The expected productive value of an applicant is $\mu_{ab}$, the taste cost attached to cue $a$ is $\tau a$, and the employer calls back when
\[
 \mu_{ab}-\tau a-k\ge\epsilon,
 \qquad \Prb(\epsilon\le z)=F(z),                    \tag{2}
\]
where $F$ is known and strictly increasing and $k$ is a common threshold. Assume the same taste cost and noise law in all regimes. For interior probabilities define the transformed difference
\[
 d_b=F^{-1}(p_{1b})-F^{-1}(p_{0b}).
\]
Equation (2) then yields
\[
 d_b=\delta_b-\tau,\qquad \delta_b=\mu_{1b}-\mu_{0b}. \tag{3}
\]
This is an index-scale statement: raw callback differences generally cannot be additively partitioned in the same way.

\begin{proposition}[Beliefs and tastes are observationally equivalent]
In an ordinary regime $b=0$, if belief levels are unrestricted, every scalar $t$ is compatible with the observed callback probabilities through $\tau=t$, $\mu_{00}=k+F^{-1}(p_{00})$ and $\mu_{10}=k+F^{-1}(p_{10})+t$. Therefore only $\delta_0-\tau$ is identified by (2).
\end{proposition}
\begin{proof}
Substitute these values in (2). The cue-zero index is $F^{-1}(p_{00})$ and the cue-one index is $F^{-1}(p_{10})$ regardless of $t$. Both callback probabilities and the entire binary conditional response law are unchanged. If tastes or beliefs have known bounds, intersect this line with those bounds; point identification still requires that the intersection have one point.
\end{proof}

For a numerical index check, suppose $d_0=-0.4$. A pure-taste model $(\tau,\delta_0)=(0.4,0)$ and a pure-belief model $(0,-0.4)$ have identical transformed callback gaps. Levels can be chosen by the proposition so the whole observed callback law agrees. The models react differently to a policy that makes productive quality perfectly observable. This is a policy-relevant nonidentification result, not a statement that either mechanism is absent.

\section{When an information factorial identifies the decomposition}
Suppose regime $b=1$ reveals validated productive quality, employers understand and trust the revelation, and the distribution of revealed quality conditional on standardized $x$ is the same under either randomized cue. Impose the structural implication $\mu_{11}=\mu_{01}$, while keeping the taste cost, threshold and index-noise law unchanged across regimes. Then (3) gives the explicit solution
\[
 \tau=-d_1,\qquad \delta_0=d_0-d_1.                  \tag{4}
\]
The proof is substitution of $\delta_1=0$ followed by the ordinary-regime equation. For example $d_0=-0.4,d_1=-0.1$ implies taste cost $0.1$ and ordinary belief difference $-0.3$ in this model. If $d_1>0$ but maintained tastes require $\tau\ge0$, the restrictions are rejected rather than repaired by silently changing the interpretation.

The identifying content comes from the equal-beliefs and stability assumptions, not merely from adding another treatment arm. A certificate that is distrusted differently by cue need not make $\delta_1=0$. Quality revelation may change attention, scrutiny or social preferences, so a regime-specific taste $\tau_1$ cannot identify $\tau_0$ through (4). An unknown regime-specific noise scale also destroys the common index comparison. Belief elicitation offers additional observables only with a stated link from incentivized reports to the beliefs used in hiring. Blind evaluation can remove both a taste trigger and predictive information; its effect is not mechanically a taste estimate.

Networks are absent from (2). To identify network effects one would need a separate intervention on referral access or exposure and a model of how it alters the opportunity set. Reinterpreting $\delta_b$ as a network share would change the model after observing the data. The full catalogue set $\mathcal T(P)$ must retain such observationally compatible alternatives unless excluded by declared measurements or restrictions.

\section{A selection reversal in career data}
Consider a fixed baseline cohort with equal numbers of high- and low-productivity individuals. Their earnings if hired are respectively $2$ and $1$; nonhire earnings are zero. Policy $\pi_0$ hires everyone. Policy $\pi_1$ hires only the high-productivity half, with no change in any person's wage conditional on hiring. Then
\[
 \E[Y\mid\text{hired},\pi_0]=1.5,\qquad
 \E[Y\mid\text{hired},\pi_1]=2,
\]
but the unconditional cohort means are
\[
 \E Y(\pi_0)=1.5,\qquad \E Y(\pi_1)=1,
 \qquad \chi=-0.5.                                  \tag{5}
\]
Thus an increase of $0.5$ in selected workers' average earnings coexists with an equally sized decrease in the correct cohort outcome. This example uses no wage discrimination at all: the reversal arises solely from selection into employment. Discounting a constant repetition of these outcomes multiplies both effects by the same positive annuity factor and preserves the reversal.

Moreover, callback-only data contain no restriction on future promotion, hours, separation or earnings beyond whatever structure is explicitly imposed. Two models can agree on every $p_{ab}(x)$ while assigning different future earnings under the career policy, producing different $\chi$. Following only successful applicants does not resolve this missing information. A randomized employer policy with baseline-cohort follow-up identifies its total career effect under assignment, interference and attrition restrictions; component effects still need feasible factorial interventions.

The remaining work is empirical and structural: specify credible quality signals, obtain comparable belief and career measurements, test the stability assumptions, and characterize the larger identified set allowing networks and accumulated disadvantage. The exact submodel inversion (4) and reversal (5) delimit what such data must add. They do not settle the multi-mechanism catalogue question.

\section{Completion Estimate}
50\%

This is a post hoc, heuristic estimate for the full catalogue target.
The attempt identifies standardized callback contrasts, proves taste-belief observational equivalence, and derives a quality-revelation inversion under strict stability assumptions. Full mechanism and career policy sets still require validated cue interpretation, belief and network interventions, accumulated-history models, and unconditional baseline-cohort earnings follow-up.

\begin{thebibliography}{9}
\bibitem{bm} M. Bertrand and S. Mullainathan (2004), \emph{Are Emily and Greg More Employable Than Lakisha and Jamal? A Field Experiment on Labor Market Discrimination}, American Economic Review 94(4), 991--1013. \href{https://www.aeaweb.org/articles?id=10.1257/0002828042002561}{Primary publisher record}. The cited experiment motivates the standardized callback target; the structural equivalence and numerical examples above are self-contained illustrations.
\end{thebibliography}

\end{document}
